\chapter{Exceptions}
\label{ch:exceptions}

\standfirst{Smalltalk-80 has no exceptions. This system has them, they are
written in Smalltalk, and the mechanism underneath is worth one paragraph
because it explains the shape of everything else.}

\section{What 1983 had instead}

\ct{Object>>error:}, which opens a debugger, and the convention of passing a
block for the failure case---\ct{remove:ifAbsent:},
\ct{detect:ifNone:}, \ct{at:ifAbsent:}. Neither can be caught by a caller. The
second is still good design and still all over the library; the first is what
exceptions replace.

\section{Catching}

\begin{st}
[1/0] on: ZeroDivide do: [:e | 'caught']            "'caught'"

[riskyThing]
    on: Error
    do: [:e | Transcript show: e messageText; cr. nil]
\end{st}

\ct{on:do:} is sent to a block, takes an exception \emph{class} and a handler
block, and answers either the protected block's value or the handler's. The
handler block may take one argument, the exception object, or none.

The class you name catches itself and all its subclasses, so \ct{on: Error}
catches \ct{ZeroDivide}, \ct{MessageNotUnderstood} and anything you define
under \ct{Error}.

\section{Signalling}

\begin{st}
Error signal                              "the class-side shortcut"
Error signal: 'something went wrong'
Error new signal: 'the same thing'
self error: 'the 1983 spelling, now an Error'
\end{st}

\ct{Object>>error:} has been rewritten to signal an \ct{Error}, so every one of
the hundreds of 1983 library methods that calls it now raises something you
can catch, and an unhandled one still opens the same debugger it always did.

\section{What the handler can do}

The exception object arrives with a set of decisions on it. Which one you use
is the interesting part of exception handling; catching is the easy part.

\begin{center}
\small
\begin{tabular}{@{}lp{0.60\textwidth}@{}}
\toprule
\ctp{e return: x} & abandon the protected block; \ctp{on:do:} answers
  \ctp{x}\\
\ctp{e return} & the same, answering \ctp{nil}\\
\ctp{e retry} & run the protected block again, from the beginning\\
\ctp{e retryUsing: aBlock} & run a different block instead\\
\ctp{e resume: x} & carry on from where \ctp{signal} was sent, which answers
  \ctp{x}. Only for resumable exceptions\\
\ctp{e pass} & decline: look for an outer handler\\
\ctp{e signal} & re-signal this exception\\
\bottomrule
\end{tabular}
\end{center}

Falling off the end of the handler block is the same as \ct{return:} with the
block's value, which is why the simple case reads so plainly.

\begin{st}
"retry with a fallback"
| attempts |
attempts := 0.
^[attempts := attempts + 1. self fetchFromNetwork]
    on: Error
    do: [:e | attempts < 3 ifTrue: [e retry] ifFalse: [e return: #cached]]
\end{st}

\subsection{Resumable and not}

\ct{Warning} is resumable: \ct{e resume: 9} makes the \ct{Warning signal:}
expression answer 9 and execution continues from there. \ct{Error} is not:
asking to resume one is itself an error. \ct{MessageNotUnderstood} \emph{is}
resumable, which is what lets a handler supply a value for a message the
receiver could not handle.

\section{Reading the exception}

\begin{center}
\small
\begin{tabular}{@{}lp{0.58\textwidth}@{}}
\toprule
\ctp{e messageText} & what was passed to \ctp{signal:}, or the class's
  \ctp{description} when nothing was\\
\ctp{e description} & a fixed sentence for the class ---
  \ct|'An error has occurred'| for \ctp{Error}. It does \emph{not} answer the
  message text\\
\ctp{e class} & the exception's class\\
\ctp{e signalContext} & the frame that signalled\\
\ctp{e message} & on \ctp{MessageNotUnderstood}, the message that failed\\
\ctp{e receiver} & on \ctp{MessageNotUnderstood}, who it was sent to\\
\bottomrule
\end{tabular}
\end{center}

\begin{stwarn}
\ct{description} is the \emph{class's} sentence, not the text you signalled
with:

\begin{st}
[Error signal: 'a'] on: Error do: [:e | e description]   "'An error has occurred'"
[Error signal: 'a'] on: Error do: [:e | e messageText]   "'a'"
\end{st}

You almost always want \ct{messageText}, which falls back to \ct{description}
when there is no text.
\end{stwarn}

\section{Cleanup}

\begin{st}
[ self doTheWork ] ensure: [ self tidyUp ]
[ self doTheWork ] ifCurtailed: [ self undoTheDamage ]
\end{st}

\ct{ensure:} runs its block on every exit: normal, non-local return, or an
exception unwinding past. \ct{ifCurtailed:} runs only on the abnormal ones.
These are the reason the concurrency classes in Part V can promise that a lock
is released.

\section{Defining your own}

Subclass \ct{Error} (or \ct{Exception}, or \ct{Warning} if it should be
resumable), and give it whatever state the handler will need:

\begin{st}
Error subclass: #InsufficientFunds
    instanceVariableNames: 'requested available'
    classVariableNames: ''
    poolDictionaries: ''
    category: 'MyStuff-Banking'

InsufficientFunds >> description
    ^'not enough money'

InsufficientFunds >> shortfall
    ^requested - available
\end{st}

\ldots and signal it with the state filled in:

\begin{st}
withdraw: anAmount
    anAmount > balance ifTrue: [
        ^(InsufficientFunds new)
            requested: anAmount available: balance;
            signal].
    balance := balance - anAmount
\end{st}

The handler can then act on the specifics rather than parsing a string, which
is the entire argument for exception \emph{classes} over error codes.

\section{What is in the box}

\begin{plain}
Exception
    Error
        ZeroDivide
        MessageNotUnderstood
        AssertionFailure
        NotFound
        SubscriptOutOfBounds
        NonBooleanReceiver
        SyntaxErrorNotification
        OutOfMemory
        TestFailure          (SUnit)
    Warning
    TimedOut
\end{plain}

The last two are raised by the VM and by the compiler rather than by library
code, and both replaced something that could not be caught at all.

\ct{NonBooleanReceiver} is what \ct|x ifTrue: [...] ifFalse: [...]| raises
when \ct{x} is neither \ct{true} nor \ct{false}. The interpreter used to
print a line and stop---no send, so nothing to catch, and the process was
simply gone along with every other worker's work. It arrives almost always
from an instance variable that \ct{new} never set, so \ct{e object} answers
the value and the message text names it. \ct{Object>>mustBeBoolean} answers
\ct{true} when nobody handles it, which is 1983's `proceed for truth', and
the interpreter retests what it answers.

\ct{SyntaxErrorNotification} is what \ct{Compiler evaluate:} raises when the
text will not parse and there is no editor to put the complaint into.
Headless, 1983 opened a window that nothing could dismiss and suspended the
process for good; \ct{on: Error do:} caught neither, because nothing was
signalled. \ct{e errorPosition} is where, and \ct{e errorCode} is the source
with the complaint spliced in at that point. An unhandled one still opens
the 1983 notifier when there is a screen for it.

\ct{OutOfMemory} is raised when there was no room to enter a method. Catching
it is what lets one worker's runaway stop being every worker's: a REST
request that exhausts the object table fails, and the server goes on serving.
Raising it costs objects at the moment there are none, so the object memory
keeps an emergency reserve and releases it once for the signal---which means
a program that catches this and then allocates on regardless does eventually
stop. Very deep stacks are the exception: past sixteen thousand frames the
interpreter stops rather than signals, because delivering an exception is
quadratic in the depth and a runaway recursion is millions of frames down.
That case still prints a line and stops, in about four seconds.

\section{How it works, in one paragraph}

Because it explains the constraints. \ct{BlockClosure>>on:do:} declares
\ct{<primitive: 199>}, and nothing implements primitive 199---an unimplemented
primitive fails, the Smalltalk body runs, and the number is left behind as a
\emph{label} on that stack frame. Signalling walks out along the sender chain
looking for a frame whose method says 199 and whose exception class matches,
runs the handler there, and answers its value from that frame. \ct{ensure:}
uses 198 the same way. Everything except the last step---returning a value
from a frame further out---is ordinary Smalltalk you can read in
\ct{lib/Kernel-Exceptions/}.

\begin{stnote}
This is Squeak's design, and it was checked against the 1983 image before
being adopted: the highest primitive that image uses is 135, so 198 and 199
were free to become labels.
\end{stnote}
